% Tutorial 1: Mt. Gox Collapse - ECOM215
\documentclass[11pt]{article}
\usepackage{fullpage}
\usepackage[left=1in,top=1in,right=1in,bottom=1in,headheight=3ex,headsep=3ex]{geometry}
\usepackage{graphicx}
\usepackage{float}
\usepackage{booktabs}
\usepackage{enumitem}

\usepackage[sc]{mathpazo}
\linespread{1.05}
\usepackage[T1]{fontenc}

\usepackage{fancyhdr,lastpage}
\pagestyle{fancy}
\usepackage{hyperref}
\usepackage{lastpage}

\usepackage{array, xcolor}
\usepackage{colortbl}
\definecolor{clemsonorange}{HTML}{EA6A20}
\definecolor{ImperialBlue}{RGB}{0, 62, 116}
\definecolor{AccentGreen}{RGB}{46, 139, 87}
\hypersetup{colorlinks,breaklinks,linkcolor=clemsonorange,urlcolor=clemsonorange,anchorcolor=clemsonorange,citecolor=black}

\lhead{}
\chead{}
\rhead{\footnotesize ECOM215: Tutorial 1 --- Mt.\ Gox Collapse}
\lfoot{}
\cfoot{\small \thepage/\pageref*{LastPage}}
\rfoot{}

\title{Tutorial 1: The Mt.\ Gox Collapse\\[0.3em] \Large Not Your Keys, Not Your Coins\\[1em] \normalsize ECOM215: Blockchain Economics and Digital Assets}
\author{Week 2 $\vert$ Foundations of Blockchain Technology}
\date{Semester B, 2025/2026}

\begin{document}

\maketitle

\vspace{1em}
\begin{center}
\fbox{\parbox{0.9\textwidth}{\centering\textbf{CASE BRIEF FOR STUDENTS}\\[0.3em] Please read before the tutorial. Estimated reading time: 15 minutes.}}
\end{center}

\section*{Background}

In early 2014, Mt.\ Gox was the world's largest Bitcoin exchange, handling approximately 70\% of all Bitcoin transactions globally. Based in Tokyo and run by French developer Mark Karpelès, the exchange had grown from a small trading card website into the dominant gateway between traditional currency and Bitcoin.

On 7 February 2014, Mt.\ Gox halted all Bitcoin withdrawals, citing ``technical issues.'' Three weeks later, the company filed for bankruptcy, announcing that 850,000 Bitcoin---worth approximately \$450 million at the time---had been stolen. At today's prices, that amount would be worth tens of billions of dollars.

The Mt.\ Gox collapse remains the largest loss in cryptocurrency history and offers fundamental lessons about the gap between blockchain's theoretical properties and how people actually use cryptocurrency.

\section*{What Was Mt.\ Gox?}

Mt.\ Gox began in 2010 as ``Magic: The Gathering Online eXchange''---a platform for trading collectible cards. Its founder, Jed McCaleb, repurposed it for Bitcoin trading and sold it to Mark Karpelès in 2011.

\medskip
\noindent\textbf{How the exchange worked:}
\begin{itemize}[leftmargin=*]
\item Users deposited fiat currency (dollars, euros, yen) or Bitcoin into Mt.\ Gox accounts
\item Mt.\ Gox held these assets on users' behalf
\item Users traded within the exchange's internal ledger
\item Users could withdraw to external bank accounts or Bitcoin wallets
\end{itemize}

\medskip
\noindent\textbf{The critical point:} When users deposited Bitcoin to Mt.\ Gox, the exchange took custody of the private keys. Users received an IOU---a balance on Mt.\ Gox's internal database---not actual control of Bitcoin on the blockchain.

This is the difference between \textit{custodial} and \textit{non-custodial} arrangements:
\begin{itemize}[leftmargin=*]
\item \textbf{Custodial}: A third party (the exchange) holds your private keys
\item \textbf{Non-custodial}: You hold your own private keys in a personal wallet
\end{itemize}

\section*{What Went Wrong}

The full picture of Mt.\ Gox's failure emerged slowly through bankruptcy proceedings and investigations. Multiple factors contributed:

\medskip
\noindent\textbf{1. Ongoing theft (undetected for years)}
\begin{itemize}[leftmargin=*]
\item Hackers had been draining Bitcoin from Mt.\ Gox's wallets since at least 2011
\item The exchange's systems failed to detect the discrepancy between customer balances and actual holdings
\item By 2014, the exchange was insolvent---it owed customers far more Bitcoin than it actually held
\end{itemize}

\medskip
\noindent\textbf{2. Poor security practices}
\begin{itemize}[leftmargin=*]
\item Private keys were stored in ``hot wallets'' (connected to the internet) rather than secure cold storage
\item Inadequate monitoring of outgoing transactions
\item No regular audits reconciling customer balances with actual reserves
\end{itemize}

\medskip
\noindent\textbf{3. Technical vulnerabilities}
\begin{itemize}[leftmargin=*]
\item Transaction malleability: Attackers could modify transaction IDs, causing Mt.\ Gox's systems to re-send payments it thought had failed
\item The exchange's code was reportedly poorly written and difficult to maintain
\end{itemize}

\medskip
\noindent\textbf{4. Management failures}
\begin{itemize}[leftmargin=*]
\item The company grew far beyond the founder's technical and operational capacity
\item Warning signs were ignored or explained away
\item No proper financial controls or governance structures
\end{itemize}

\section*{The Irony}

Bitcoin was designed to eliminate the need for trusted intermediaries. Its core properties---cryptographic security, transparency, immutability---allow peer-to-peer transactions without banks or payment processors.

Yet the largest Bitcoin disaster occurred not because of any failure in the Bitcoin protocol, but because users \textbf{reintroduced a trusted intermediary}. They handed their private keys to an exchange, trusting Mt.\ Gox to safeguard their assets.

\medskip
\noindent\textbf{The blockchain worked exactly as designed:}
\begin{itemize}[leftmargin=*]
\item Every Bitcoin transaction was valid and properly recorded
\item The cryptography was never broken
\item The Bitcoin network continued operating without interruption
\end{itemize}

\medskip
\noindent\textbf{The failure was at the interface:}
\begin{itemize}[leftmargin=*]
\item Users trusted a centralised company instead of using the decentralised protocol directly
\item The exchange held private keys, creating a single point of failure
\item When the exchange failed, users had no recourse on the blockchain---their Bitcoin was gone
\end{itemize}

This pattern---blockchain technology working correctly while intermediaries built on top of it fail---has repeated throughout cryptocurrency history.

\section*{The Aftermath}

\textbf{For Mt.\ Gox users:}
\begin{itemize}[leftmargin=*]
\item Approximately 24,000 creditors lost funds
\item Bankruptcy proceedings lasted over a decade
\item A portion of Bitcoin was later recovered (200,000 BTC found in an old wallet)
\item Creditors began receiving partial repayments in 2024---ten years after the collapse
\end{itemize}

\medskip
\noindent\textbf{For Mark Karpelès:}
\begin{itemize}[leftmargin=*]
\item Arrested in Japan in 2015
\item Charged with embezzlement and data manipulation
\item Found guilty of falsifying records but acquitted of embezzlement
\item Received a suspended sentence
\end{itemize}

\medskip
\noindent\textbf{For the industry:}
\begin{itemize}[leftmargin=*]
\item The phrase ``not your keys, not your coins'' became a mantra
\item Proof-of-reserves became a standard expectation for exchanges
\item Hardware wallets and cold storage solutions gained popularity
\item Regulatory attention increased, eventually leading to frameworks like MiCA
\end{itemize}

\section*{Key Facts Summary}

\begin{center}
\begin{tabular}{ll}
\toprule
\textbf{Item} & \textbf{Detail} \\
\midrule
Peak market share & $\sim$70\% of global Bitcoin trading \\
Date of collapse & February 2014 \\
Bitcoin lost & 850,000 BTC (650,000 customer + 200,000 company) \\
Value at time & $\sim$\$450 million \\
Value at 2024 prices & $\sim$\$50+ billion \\
Cause & Ongoing theft + poor security + management failures \\
Resolution & Partial repayment began 2024 (10 years later) \\
\bottomrule
\end{tabular}
\end{center}

\section*{Questions to Consider}

\begin{enumerate}[leftmargin=*]
\item Bitcoin was designed to eliminate trusted intermediaries. Why did users hand their Bitcoin to Mt.\ Gox anyway?

\item The blockchain worked perfectly throughout the Mt.\ Gox collapse. Does this mean the technology succeeded or failed?

\item ``Not your keys, not your coins'' is good advice. But is it realistic to expect ordinary users to manage their own private keys?

\item How should we think about the trade-off between the security of self-custody and the convenience of using intermediaries?

\item Mt.\ Gox would likely be illegal under today's regulations (MiCA, etc.). Does regulation solve the problem, or just shift the trust to regulators?
\end{enumerate}

\section*{Further Reading (Optional)}

\begin{itemize}[leftmargin=*]
\item ``The Mt.\ Gox Bitcoin Debacle'' - Wired (2014)
\item Kim Nilsson's investigation: ``The Missing MtGox Bitcoins'' (WizSec blog)
\item Mt.\ Gox bankruptcy trustee reports (available online)
\end{itemize}

\newpage

%==============================================================================
% PART B: TA DISCUSSION GUIDE
%==============================================================================

\section*{Session Timeline}

\begin{center}
\begin{tabular}{ll}
\toprule
\textbf{Time} & \textbf{Activity} \\
\midrule
0:00--0:08 & Context setting: What Mt.\ Gox was and what happened \\
0:08--0:22 & Discussion Question 1: Why intermediaries persist \\
0:22--0:36 & Discussion Question 2: Did blockchain succeed or fail? \\
0:36--0:48 & Discussion Question 3: Self-custody vs.\ convenience \\
0:48--0:55 & Discussion Question 4: Role of regulation \\
0:55--1:00 & Synthesis and key takeaways \\
\bottomrule
\end{tabular}
\end{center}

\section*{Discussion Questions with Guidance}

\subsection*{Question 1: Why did users trust Mt.\ Gox?}

\textit{``Bitcoin was created specifically to eliminate the need for trusted intermediaries. Yet users voluntarily handed their Bitcoin to a centralised exchange. Why?''}

\medskip
\noindent\textbf{Points that may emerge:}

\begin{itemize}[leftmargin=*]
\item \textbf{Convenience}: Managing private keys is technically demanding. Exchanges offer a familiar banking-like experience.
\item \textbf{Trading functionality}: To actively trade, you need assets on an exchange. The blockchain itself doesn't provide order books or instant execution.
\item \textbf{Fiat on/off-ramps}: Converting between dollars/euros and Bitcoin requires an intermediary in most jurisdictions.
\item \textbf{Network effects}: Mt.\ Gox had the most liquidity, so traders gravitated there.
\item \textbf{Lack of alternatives}: In 2011--2014, user-friendly wallets and hardware solutions were primitive.
\item \textbf{Complacency}: Mt.\ Gox seemed established and reliable. ``Too big to fail'' mentality.
\end{itemize}

\medskip
\noindent\textbf{Key insight}: Blockchain technology provides the \textit{option} of trustless transactions, but doesn't \textit{force} users to take it. In practice, convenience often wins over security. This creates what we might call the \textbf{intermediary re-emergence problem}---new intermediaries form around decentralised systems.

\subsection*{Question 2: Did blockchain succeed or fail?}

\textit{``The Bitcoin blockchain worked perfectly throughout Mt.\ Gox's collapse. Every transaction was valid, the network kept running, no cryptography was broken. So did the technology succeed---or did the overall system fail?''}

\medskip
\noindent\underline{Arguments that blockchain succeeded:}
\begin{itemize}[leftmargin=*]
\item The protocol did exactly what it was designed to do
\item Users who held their own keys were unaffected
\item The theft was detectable precisely because of blockchain transparency
\item Mt.\ Gox was not ``the blockchain''---it was a company built on top of it
\end{itemize}

\medskip
\noindent\underline{Arguments that the system failed:}
\begin{itemize}[leftmargin=*]
\item The \textit{ecosystem} failed, even if the protocol didn't
\item Technology must be evaluated by real-world outcomes, not just theoretical properties
\item If most users need intermediaries to use blockchain, then those intermediaries are part of the system
\item \$450 million lost is a failure by any practical measure
\end{itemize}

\medskip
\noindent\textbf{Key insight}: This question highlights the distinction between \textit{protocol-level} properties and \textit{system-level} outcomes. Blockchain provides certain guarantees \textit{if used as designed}. But how people actually use it determines real-world results. A system that most users cannot safely use directly has practical limitations.

\subsection*{Question 3: Self-custody vs.\ convenience}

\textit{``Not your keys, not your coins'' is sound advice. But managing private keys---backing them up, securing them, never losing them---is demanding. Is self-custody realistic for ordinary users? What's the right balance?''}

\medskip
\noindent\underline{Case for self-custody:}
\begin{itemize}[leftmargin=*]
\item Mt.\ Gox, FTX, Celsius, Voyager---the list of failed custodians is long
\item Hardware wallets (Ledger, Trezor) have made self-custody much easier
\item ``Be your own bank'' is blockchain's core value proposition
\item Long-term holders who don't need to trade have less need for exchanges
\end{itemize}

\medskip
\noindent\underline{Case for custodial solutions:}
\begin{itemize}[leftmargin=*]
\item Most people struggle with password management; private keys are harder
\item Lost keys mean permanently lost assets---no recovery, no customer service
\item Regulated custodians (Coinbase Custody, Fidelity) now offer institutional-grade security
\item Active traders need exchange access; moving on/off chain for every trade is impractical
\end{itemize}

\medskip
\noindent\underline{Middle ground options:}
\begin{itemize}[leftmargin=*]
\item Multi-signature wallets (require multiple keys to transact)
\item Social recovery (trusted contacts can help restore access)
\item Keep trading amounts on exchanges, long-term holdings in self-custody
\end{itemize}

\medskip
\noindent\textbf{Key insight}: There's a genuine trade-off between \textit{trustlessness} and \textit{usability}. Maximising one often means sacrificing the other. The right choice depends on the user's technical sophistication, risk tolerance, and use case. This is not a solved problem.

\subsection*{Question 4: Does regulation solve the problem?}

\textit{``Under MiCA and similar regulations, exchanges must maintain reserves, segregate customer assets, and submit to audits. Mt.\ Gox would be illegal today. Does regulation solve the custody problem?''}

\medskip
\noindent\underline{Arguments that regulation helps:}
\begin{itemize}[leftmargin=*]
\item Minimum standards prevent the worst abuses
\item Proof-of-reserves and audits create accountability
\item Regulated entities can be sued, prosecuted, shut down
\item Ordinary users can't audit exchanges themselves---regulators can
\end{itemize}

\medskip
\noindent\underline{Arguments that regulation has limits:}
\begin{itemize}[leftmargin=*]
\item FTX was regulated in multiple jurisdictions and still collapsed
\item Regulation shifts trust from the exchange to the regulator---it doesn't eliminate trust
\item Enforcement is reactive, not preventive; fraud is usually discovered after the damage
\item Global market: Users can access unregulated offshore exchanges
\end{itemize}

\medskip
\noindent\underline{The deeper tension:}
\begin{itemize}[leftmargin=*]
\item Blockchain was designed to eliminate the need for trusted institutions
\item Regulation reintroduces institutional oversight
\item This may be necessary and practical, but it changes the nature of what blockchain offers
\end{itemize}

\medskip
\noindent\textbf{Key insight}: Regulation doesn't eliminate the trust problem---it redirects it. Instead of trusting the exchange alone, you trust the regulatory system (laws, enforcement, auditors). This may be a reasonable trade-off, but it's different from the trustless vision of blockchain's original design.

\subsection*{Extension Question (if time permits)}

\textit{``In 2024, spot Bitcoin ETFs were approved in the US. These are custodial products---investors don't hold private keys; they hold shares in a fund that holds Bitcoin. Is this a betrayal of Bitcoin's principles, or a pragmatic path to adoption?''}

\medskip
This connects Week 1 material to the later topic on traditional finance applications. ETFs bring institutional custody, regulatory oversight, and mainstream access---but they also reintroduce all the intermediaries Bitcoin was designed to bypass.

\vfill
\begin{flushright}
\textit{End of Tutorial 1 Materials}
\end{flushright}

\end{document}